%% \documentclass{article}

%% \usepackage[english]{babel}

%% \usepackage[a4paper,top=2cm,bottom=3cm,left=3cm,right=2cm,marginparwidth=1.75cm]{geometry}

%% % Useful packages
%% \usepackage{amsmath}
%% \usepackage{amssymb}
%% \usepackage{graphicx}
%% \usepackage[colorlinks=true, allcolors=blue]{hyperref}
%% \usepackage{tikz}
%% \usepackage{float}
%% \usepackage{graphicx}
%% \usepackage{subcaption}
%% \usepackage{todonotes}
%% \usepackage{arydshln}
%% \usepackage{mathtools}
%% \usepackage{bussproofs}
%% \usepackage{xfrac}
%% \usepackage{changepage}
%% \usepackage{todonotes}
%% \usepackage{semantic}
%% \usepackage{amsthm}
%% \usepackage{stmaryrd}
%% \usepackage{algorithm}
%% \usepackage[noend]{algpseudocode}
%% \newcommand{\impl}{\Rightarrow }
%% \renewcommand{\iff}{\Leftrightarrow }
%% \renewcommand{\phi}{\varphi}
%% \renewcommand{\epsilon}{\varepsilon}
%% %% \newcommand{\land}{\wedge }
%% %% \newcommand{\lor}{\vee }
%% \newcommand{\A}{\mathcal{A} }
%% \newcommand{\D}{\mathcal{D} }
%% \newcommand{\I}{\mathcal{I} }
%% \renewcommand{\P}{\mathcal{P} }
%% \newcommand{\F}{\mathcal{F} }

%% \tikzset{node/.style={circle,text=black, minimum size=.8cm, draw}}

%% \newtheorem{theorem}{Theorem}[section]
%% \newtheorem{definition}[theorem]{Definition}
%% \newtheorem{lemma}[theorem]{Lemma}
%% \newtheorem{corollary}[theorem]{Corollary}
%% \newtheorem{proposition}[theorem]{Proposition}
%% \newtheorem{remark}[theorem]{Remark}
%% %\theoremstyle{definition}
%% \newtheorem{exmp}[theorem]{Example}

%% \title{Assignment 10}
%% \author{Alex Loitzl}

%% \begin{document}
%% \maketitle

\section{Homework 9}

\subsection{Formulate QBF as a first-order theory}
We define the theory $T_{QBF}$ over the signature $\langle \{t^0, f^0\}, \{v^1\} \}\rangle$ with constant function symbols $t$ and $f$ and a unary predicate $v$.

\subsubsection{Define a first-order theory $T_{QBF}$}
We axiomatize the theory:
\begin{enumerate}
  \item $\forall x, x = t \lor x = f$
  \item $t \neq f$
  \item $v(t)$
  \item $\neg v(f)$
\end{enumerate}

The first two axioms restrict the models to exactly those models with a domain of cardinality of two. The third and fourth axiom ensure that the constant $t$ is interpreted as true and $f$ as false.

\subsubsection{Define a function f from closed QBF formulas to $T_{QBF}$ formulas s.t. a formula $\phi$ is true iff $f(\phi)$ is satisfiable.}

Let $\mathcal{P}$ be the set of propositions. We define a function $f : \mathcal{P} \to \mathcal{A}$ from propositions to atomic formulas.
We consider formulas as defined in the lecture and add existential quantification and $\top$ as they are used in QBF.
\begin{itemize}
  \item $f(\top) \mapsto \top$
  \item $f(\bot) \mapsto \bot$
  \item $f(p) \mapsto v(p)$
\end{itemize}
We then get a function $\mathcal{F}: \psi_{QBF} \to \psi_{T_{QBF}}$ by mapping $f$ over the connectives of FOL.

\begin{lemma}
  For any closed QBF formula $\phi$, $\phi$ is true $\iff$ $\mathcal{F}(\phi)$ is satisfiable.
\end{lemma}
\begin{proof}$ $
  We will consider the natural structure $\mathcal{I} \coloneqq \langle \{0, 1\}, \{t,f\}, \{v\}\rangle$ with $t_{\mathcal{I}} = 1$, $f_\mathcal{I} = 0$ and $v = \{1\}$. We omit writing the subscript $\mathcal{I}$ in the following. Clearly, the structure is a $T_{QBF}$-$Model$. We prove that $\phi$ is true \emph{iff} $\mathcal{I} \models \mathcal{F}(\phi)$ by induction on the structure of $\phi$.
  \begin{itemize}
    \item $\phi \coloneqq \bot$. For both directions we get an immediate contradiction as $\mathcal{I} \not\models \bot$ and $\bot$ is not true.
    \item $\phi \coloneqq \top$, trivially $\mathcal{I} \models \mathcal{F}(\top) = \top$ and $\top$ is true, hence $\top$ is true \emph{iff} $\mathcal{I} \models \mathcal{F}(\top)$.
    \item $\phi \coloneqq p$, (for completenes, $\phi$ is closed)
    \item $\phi \coloneqq \phi_1 \impl \phi_2$. Assume $\phi$ is true, then either $\phi_1$ is false or both $\phi_1$ and $\phi_2$ are true. By induction, in the first case we get $\mathcal{I} \not\models \mathcal{F}(\phi_1)$, and therefore $\mathcal{I} \models \mathcal{F}(\phi)$ and in the second case $\mathcal{I} \models \mathcal{F}(\phi_1)$ and $\mathcal{I} \models \mathcal{F}(\phi_2)$, hence $\mathcal{I} \models \mathcal{F}(\phi)$. For the other direction, assume $\mathcal{I} \models \mathcal{F}(\phi)$. Then either $\mathcal{I} \not\models \mathcal{F}(\phi_1)$ or both $\mathcal{I} \models \mathcal{F}(\phi_1)$ and $\mathcal{I} \models \mathcal{F}(\phi_2)$. By induction, in the first case $\phi_1$ is false, hence $\phi$ true and in the second case both $\phi_1$ and $\phi_2$ are true, hence $\phi$ is true.
    \item $\phi \coloneqq \forall \phi_1$. By (IH) we know that $\phi_1[\top]$ is true \emph{iff} $\mathcal{I} \models \mathcal{F}(\phi_1[\top])$ and $\phi_1[\bot]$ is true \emph{iff} $\mathcal{I} \models \mathcal{F}(\phi_1[\bot])$.
    We continute to show that $\mathcal{I} \models \mathcal{F}(\phi_1[\top])$ \emph{iff} $\mathcal{I}_{x \gets 1} \models \mathcal{F}(\phi_1)$ and $\mathcal{I} \models \mathcal{F}(\phi_1[\bot])$ \emph{iff} $\mathcal{I}_{x \gets 0} \models \mathcal{F}(\phi_1)$, hence by definition of $\forall$, $\mathcal{I} \models \mathcal{F}(\phi)$
      \paragraph{case $\top$:} Since $\mathcal{I} \models \top \iff v(t)$, we know $\mathcal{I} \models \mathcal{F}(\phi_1[v(t)])$ \emph{iff} $\mathcal{I} \models \mathcal{F}(\phi_1[\top])$ (\emph{Replacement theorem} would require another (trivial) proof by induction on formula). Moreover, $\mathcal{I} \models \mathcal{F}(\phi_1[v(t)]) \iff \mathcal{I}_{x \gets 1} \models \mathcal{F}(\phi_1[v(x)])$ and $\phi_1[v(x)] = \phi_1$. By transitivity of $\iff$, $\phi_1[\top]$ is true \emph{iff} $\mathcal{I}_{x \gets 1} \models \mathcal{F}(\phi_1)$.
      \paragraph{case $\bot$:} Symmetric to \textbf{case $\top$}.
    \item $\phi \coloneqq \exists \phi_1$. By definition, either $\phi_1[\bot]$ or $\phi_1[\top]$ is true.
      We proceed by cases, by (IH) we know that in the first case $\phi_1[\top]$ is true \emph{iff} $\mathcal{I} \models \mathcal{F}(\phi_1[\top])$ and in the other case $\phi_1[\bot]$ is true \emph{iff} $\mathcal{I} \models \mathcal{F}(\phi_1[\bot])$. For each case, the proof is identical to the cases $\bot$ and $\top$ above.
  \end{itemize}

  To prove the final statement, assume $\phi$ is true, then by the proof above $\mathcal{F}(\phi)$ is satisfiable as $\mathcal{I} \models \mathcal{F}(\phi)$. For the other direction, assume $\mathcal{F}(\phi)$ is satisfiable, then there exists an $\mathcal{I}'$ s.t. $\mathcal{I}' \models \mathcal{F}(\phi)$. By the axioms, $\mathcal{I}'$ is isomorphic to $\mathcal{I}$, since $\mathcal{I}'$ is a $T_{QBF}$-Model and therefore by axioms (1) and (2) has a 2 element domain, with the element $x = \llbracket t \rrbracket_\mathcal{I}'$ being the only element s.t. $x \in \llbracket v\rrbracket_\mathcal{I}'$. Thereofre for any $\psi$, $\mathcal{I}' \models \psi$ \emph{iff} $\mathcal{I} \models \psi$. Hence, also $\mathcal{I} \models \mathcal{F}(\phi)$ and therefore $\phi$ is true.
\end{proof}


\subsection{Argue that Pressburger arithmetic does not allow QE.}

The goal of this exercise is to show that Presburger arithmetic $\TPR$
as defined in the lecture (i.e., without divisibility predicate) does \emph{not}
allow for quantifier elimination (QE).

It is important to point out that in most cases, Presburger arithmetic is
considered with the divisibility predicate, which allows for quantifier
elimination, this construction was even given by Presburger in his original 1929 paper.

To show that QE does not work without the divisibility predicate, consider the
$\TPR$ formula $\PH$ that expresses that a variable $y$ is even.

\[
\PH:~~\exists x.~x + x \leq y \land y \leq x + x
\]

Note that in this formula, $y$ is a \emph{free variable}. Depending on the
value of $y$, the formula $\PH$ is either true or false. The set of
values of $y$ for which $\PH$ is true is the set of even numbers, which
is countably infinite.

For the sake of readability, we say a formula $\PS$ with a single free variable $x$ \emph{defines} a subset
of the natural numbers $S$ if $\PS$ evaluates to true if and only if $x \in S$.
We now claim that any \emph{quantifier-free} $\TPR$-formula $\PS$ with a single free variable
defines a subset of the natural numbers that can be expressed as
a finite union of intervals. To show this, we first consider atomic formulas
and then proceed by induction on the structure of $\PS$.

\begin{lemma}
  \label{lem:atomic}
  Every quantifier-free, atomic $\TPR$-formula $\PS$ with a single free variable $x$
  defines a contiguous subset of the natural numbers.
\end{lemma}

\begin{proof}
  Denoting repeated addition of $x$ by $ax$ using some constant $a$, we can
  express any atomic formula in $\TPR$ as $ax + b \leq cx + d$, which simplifies
  to $(a-c)x \leq d - b$. This atomic formula is true for integer values of $x$
  in the following intervals:
  \begin{itemize}
    \item When \( a > c \):
      \begin{itemize}
        \item If \( d \geq b \), then \( x \in \left[0, \frac{d - b}{a - c}\right] \).
        \item If \( d < b \), then \( x \in \left[\frac{d - b}{a - c}, \infty\right) \). However, since \( \frac{d - b}{a - c} \) is negative and $x \in \mathbb N$, this effectively means \( x \in [0, \infty) \).
      \end{itemize}

    \item When \( a < c \):
      \begin{itemize}
        \item If \( d \geq b \), then \( x \in \left[\frac{d - b}{a - c}, \infty\right) \).
        \item If \( d < b \), then \( x \in \left[0, \frac{d - b}{a - c}\right] \).
      \end{itemize}

    \item When \( a = c \):
      \begin{itemize}
        \item If \( d \geq b \), then \( x \in \mathbb{N} \).
        \item If \( d < b \), the formula is false, i.e. $x$ is in the empty interval.
      \end{itemize}
  \end{itemize}
\end{proof}

\begin{theorem}
  Any quantifier-free $\TPR$-formula $\PS$ with a single free variable $x$ defines
  a finite union of intervals.
\end{theorem}

\begin{proof}
  We proceed by induction on the structure of $\PS$.

  \textbf{Base case}: By Lemma \ref{lem:atomic}, if $\PS$ is an atomic formula,
  it defines a contiguous subset of the natural numbers (i.e., a single
  interval).
  \textbf{Induction step}: Let $\PS$ be a boolean combination of formulas
  $\PH_1$ and $\PH_2$, that both describe finite unions of intervals.
  \begin{itemize}
    \item $\PS := \neg \PH_1$: This describes the complement of the intervals defined by $\PH_1$.
      The complement of a finite union of intervals is a finite union of intervals.
    \item $\PS := \PH_1 \land \PH_2$: This describes the intersections of the intervals defined by $\PH_1$ and $\PH_2$.
      The intersection of two finite unions of intervals is a finite union of intervals.
    \item $\PS := \PH_1 \lor \PH_2$: This describes the union of the intervals defined by $\PH_1$ and $\PH_2$.
      The union of two finite unions of intervals is still a finite union of intervals.
  \end{itemize}
\end{proof}

Observe that every finite union of intervals can only describe a subset of the natural numbers of
a certain structure:

\begin{itemize}
  \item A finite set of natural numbers, if the union does not contain the interval $\left[k, \infty\right)$ for some $k \in \mathbb Q^{+} \cup \{0\}$.
  \item Otherwise, the union of a finite set of natural numbers and $\{n \mid n \in \mathbb N \land n \geq k \}$ for some $k \in \mathbb Q^{+} \cup \{0\}$.
\end{itemize}

It is easy to see that thus, every subset of natural numbers defined by a finite
union of intervals is either \emph{finite} or \emph{has a finite complement} (when it contains
$[k, \infty)$ for some $k$).  This is in contrast to the definition of any
periodic set of natural numbers like the even numbers, which is both countably
infinite and has a countably infinite complement. Therefore, any periodic
set of natural numbers cannot be described by a quantifier-free $\TPR$-formula.
